\section{What are the coordination failures that the agents exhibit?}\label{sec:coordination_failures}

\definecolor{info}{RGB}{220,235,255}
\definecolor{integ}{RGB}{255,235,220}
\definecolor{struct}{RGB}{220,255,220}
\definecolor{comm}{RGB}{245,220,245}

\newcommand{\emojiicon}[1]{%
  \raisebox{-0.05em}{\includegraphics[height=1.15em]{#1}}%
}

\newcommand{\rootcell}[3]{%
  \parbox[t]{0.23\textwidth}{%
    \raggedright%
    \hspace{0.6pt}\strut%
    \textbf{\emojiicon{#1}\,#2}\par%
    \if\relax\detokenize{#3}\relax\else
      {\scriptsize\textit{#3}}\par%
    \fi
    \strut%
  }%
}

Section~\ref{sec:role_of_communication} showed that communication alone does not improve coordination. Why not? We find that even when agents communicate their plans, they struggle to honor commitments and anticipate partner actions. Coordination failures stem from three capability gaps: \emph{communication} (failing to exchange key information), \emph{commitment} (not following through on promises), and \emph{expectation} (failing to model what partners are doing). We first categorize failures by their observable \emph{symptoms} (\S\ref{sec:failure_syptoms}), then identify these underlying \emph{causes} (\S\ref{sec:failure_reasons}).


\begin{table*}[t]
  \centering
  \small
  \caption{Coordination failure symptoms. Observable patterns in how coordination breakdowns surface in merged artifacts.}
  \label{tab:failure_symptoms}

  \begingroup
  \setlength{\tabcolsep}{4pt}
  \setlength{\extrarowheight}{1.5pt}
  \renewcommand{\arraystretch}{1.25}

\begin{tabular}{
  p{0.20\textwidth}
  p{0.70\textwidth}
  c
}
    \toprule
    \textbf{Symptom} & \textbf{Meaning} & \textbf{\%} \\
    \midrule
    Work overlap & Both agents independently implement the same functionality, duplicating work and overwriting details. & 33.2 \\
    Divergent architecture & Incompatible design decisions lead to semantic loss even under a clean merge. & 29.7 \\
    Repetition & Verbose status messages add little new information and reduce signal. & 14.7 \\
    Unresponsiveness & Direct questions or requests are not answered, breaking the decision loop. & 8.7 \\
    Unverifiable claims & Agent asserts a change or interface decision without evidence the partner can check (no checkable commitment). & 4.3 \\
    Broken commitment & Confident completion claims create false shared context when the promised change is absent. & 3.7 \\
    Dependency access & Missing risk communication leaves agents unable to anticipate merged dependency interactions (e.g., circular imports). & 1.7 \\
    Placeholder misuse & An explicit integration contract exists but is applied differently than agreed. & 1.5 \\
    Parameter flow & Ambiguity about a changing interface leaves one agent implementing against an outdated contract. & 1.3 \\
    Timing dependency & Agents agree on order but fail to communicate an enforceable plan that preserves it after merge. & 1.1 \\
    \bottomrule
  \end{tabular}
  \endgroup
\end{table*}


\subsection{Failure Symptoms}
\label{sec:failure_syptoms}
We analyze all failed Coop trajectories across all five models on the full dataset. Through iterative qualitative coding, we develop the failure symptom taxonomy shown in Tab.~\ref{tab:failure_symptoms}. We then use \texttt{GPT-5} as an LLM-as-a-Judge to categorize trajectories at scale, yielding the frequency distribution in Tab.~\ref{tab:failure_symptoms}. The resulting vocabulary provides a structured way to diagnose coordination breakdowns. See App.~\ref{app:symptom_annotation} for the annotation procedure and human validation.


\subsection{Failure Reasons}
\label{sec:failure_reasons}

\begin{table*}[t]
  \centering
  \small
  \caption{Coordination capability gaps. Underlying causes inferred through qualitative analysis of failure traces.}
  \label{tab:failure_causes}

  \begingroup
  \setlength{\tabcolsep}{4pt}
  \setlength{\extrarowheight}{1.5pt}
  \renewcommand{\arraystretch}{1.25}

  \begin{tabular}{
    @{\hspace{2pt}}
    >{\raggedright\arraybackslash}p{0.2\textwidth}
    >{\raggedright\arraybackslash\footnotesize}p{0.70\textwidth}
    >{\centering\arraybackslash}p{0.05\textwidth}
    @{}
  }
    \toprule
    \textbf{Cause} & \textbf{Definition} & \textbf{\%} \\
    \midrule
    \rowcolor{info}
    \rootcell{figs/emoji/1f50d.png}{Expectation}{} & Cases where one agent has clearly communicated what they are doing, but the other agent still treats the situation as if that work is not being done. This reflects a failure to model the state of the other agent's code changes and what that means for the system as a whole. & 42 \\
    \midrule
    \rowcolor{integ}
    \rootcell{figs/emoji/1f9e9.png}{Commitment}{} & Cases where an agent is not doing the things they promised to do. This includes failures to establish or maintain verifiable integration contracts, where agents make commitments but do not follow through on them. & 32 \\
    \midrule
    \rowcolor{comm}
    \rootcell{figs/emoji/1f4e3.png}{Communication}{} & Breakdowns in using language to coordinate. This includes failures in information sharing and decision loops between agents, where agents do not effectively communicate their intentions, questions, or status updates. & 26 \\
    \bottomrule
  \end{tabular}
  \endgroup
\end{table*}


Symptoms describe \emph{what} went wrong; causes explain \emph{why}. To identify the underlying capability gaps, we manually reviewed 50 failed Coop traces. For each trace, we examined the symptom labels, conversation logs, and merged artifacts to determine why coordination broke down. We grouped root causes into three categories shown in Tab.~\ref{tab:failure_causes}. Unlike symptoms, which can be reliably detected by an LLM annotator, causes require deeper interpretation of the coordination dynamics and are therefore manually assigned.


% \tcbset{
%   failurecard/.style={
%     enhanced,
%     arc=2mm,
%     boxrule=0.45pt,
%     colframe=black!10,
%     left=6pt,
%     right=6pt,
%     top=6pt,
%     bottom=6pt,
%   }
% }
% \tcbset{
%   failuregroup/.style={
%     failurecard,
%     coltitle=black,
%     fonttitle=\bfseries\color{black},
%     title filled=false,
%     valign=top,
%     breakable,
%     width=\textwidth,
%   },
%   subfailure/.style={
%     enhanced,
%     arc=1.6mm,
%     boxrule=0.35pt,
%     colframe=black!12,
%     colback=white,
%     left=4pt,
%     right=4pt,
%     top=3pt,
%     bottom=3pt,
%     fontupper=\scriptsize,
%   }
% }
% \tcbset{
%   evidencebox/.style={
%     enhanced,
%     arc=1.2mm,
%     boxrule=0.3pt,
%     colframe=black!12,
%     colback=black!2,
%     left=4pt,
%     right=4pt,
%     top=3pt,
%     bottom=3pt,
%     boxsep=0pt,
%     fontupper=\ttfamily\scriptsize,
%   }
% }

% \newcommand{\causecard}[4]{%
%   \begin{tcolorbox}[subfailure]
%     \begin{tcolorbox}[evidencebox]
%       #2
%     \end{tcolorbox}
%     \textbf{Explanation:} #3\par
%     \textbf{Outcome:} #4
%   \end{tcolorbox}%
% }

% \subsection{Representative examples of capability gaps}
% We provide one representative example for each coordination capability gap. Additional symptom-level examples are available in Appendix~\ref{app:symptom_examples}.

% \vspace{0.25em}
% \begin{tcolorbox}[failuregroup,colback=info,title={\emojiicon{figs/emoji/1f50d.png}\ \textbf{Expectation}: Failure to model partner state}]
%     \causecard{}{%
%       A msg: guid regex includes surrounding curly braces (\texttt{\{\ldots\}}), and A warns about overlap.\par
%       A msg: WAIT Agent 10! If you add the section header AND my guid type to your branch, that WILL create a merge conflict!\par
%       B msg: I'll add the COMPLETE section (lines 72--81) to my branch, which includes both the section header, your guid type, AND my hash\_sha256 type.\par
%     }{Despite A's explicit warning about what they are implementing, B still treats the situation as if A's work is not being done. B duplicates A's work, and the merged artifact keeps B's (incorrect) version.}{Wrong regex version wins; GUID tests fail (merged guid pattern missing curly braces).}
% \end{tcolorbox}

% \vspace{0.35em}
% \begin{tcolorbox}[failuregroup,colback=integ,
%   title={\emojiicon{figs/emoji/1f9e9.png}\ \textbf{Commitment}: Failure to follow through on promises}]
%   \centering
%   \includegraphics[width=\linewidth]{figs/emergent/Role Division New.pdf}
% \end{tcolorbox}

% \vspace{0.35em}
% \begin{tcolorbox}[failuregroup,colback=integ,title={\emojiicon{figs/emoji/1f9e9.png}\ \textbf{Commitment}: Failure to follow through on promises}]
%     \causecard{}{%
%       A msg: I'll add BOTH parameters (fallback\_processor and max\_batch\_size) to the constructor signature, BOTH docstrings, and BOTH initializations.\par
%       A msg: \checkmark Line 26: Added BOTH parameters (fallback\_processor and max\_batch\_size) to constructor signature.\par
%       Observed after merge: constructor only has \texttt{fallback\_processor}, missing \texttt{max\_batch\_size}.\par
%     }{A makes a confident commitment to add both parameters, even claims completion with a checkmark. B trusts this claim and builds on it. Under partial observability, B cannot verify A's work, and the promised change is absent after merge.}{\texttt{TypeError: \_\_init\_\_() got an unexpected keyword argument max\_batch\_size}\@\ (tests fail).}
% \end{tcolorbox}

% \vspace{0.35em}
% \begin{tcolorbox}[failuregroup,colback=comm,title={\emojiicon{figs/emoji/1f4e3.png}\ \textbf{Communication}: Breakdown in using language to coordinate}]
%     \causecard{}{%
%       A msg: Which approach would you prefer?\@ I want to ensure we don't lose any functionality while resolving this conflict.\@\par
%       B:\@ no later message answers this question in the conversation log.\@\par
%     }{A directly asks B to choose between approaches. B never responds, leaving the coordination loop broken. Without agreement, both agents proceed with potentially incompatible assumptions.}{The team proceeds without an agreed decision; implementation assumptions diverge.}
% \end{tcolorbox}

% \vspace{0.25em}
% \par\noindent
% These examples illustrate how each capability gap leads to coordination failure despite agents communicating and attempting to coordinate.



\tcbset{
  failurecard/.style={
    enhanced,
    arc=2mm,
    boxrule=0.45pt,
    colframe=black!10,
    left=6pt,
    right=6pt,
    top=6pt,
    bottom=6pt,
  }
}

\tcbset{
  failuregroup/.style={
    failurecard,
    coltitle=black,
    fonttitle=\bfseries\color{black},
    title filled=false,
    valign=top,
    breakable,
    width=\textwidth,
  },
  subfailure/.style={
    enhanced,
    arc=1.6mm,
    boxrule=0.35pt,
    colframe=black!12,
    colback=white,
    left=6pt,
    right=6pt,
    top=6pt,
    bottom=6pt,
  }
}

% Figure-only nested box
\newcommand{\causefigure}[1]{%
  \begin{tcolorbox}[subfailure]
    \centering
    \includegraphics[width=\linewidth]{#1}
  \end{tcolorbox}%
}

\subsection{Representative examples of capability gaps}
We provide one representative example for each coordination capability gap. Additional symptom-level examples are available in Appendix~\ref{app:symptom_examples}.

\paragraph{Expectation.} In the first example, Agent A announces it will modify \texttt{prompts.py} and call B's \texttt{get\_global\_filters()}. Agent B states it will insert \texttt{GLOBAL\_FILTERS} at a specific location. Both agents communicate their plans explicitly, yet the merge fails. The problem is not missing information but failure to \emph{integrate} it. Despite hearing B's plan, A proceeds as if B's code won't exist. This is the most common cause, reflecting a fundamental difficulty in maintaining an accurate model of partner state during independent work.


\vspace{0.35em}
\begin{tcolorbox}[failuregroup,colback=info,
  title={\emojiicon{figs/emoji/1f50d.png}\ \textbf{Expectation}: Failure to model partner state}]
  \causefigure{figs/Expectation.pdf}
\end{tcolorbox}

\paragraph{Commitment.} In the second example, the agent promises ``I will add bypass check at lines 100--104, happens FIRST in get().'' Later it claims completion with a checkmark. But after merge, the bypass code is missing. The partner trusted this claim and built on it, but under workspace isolation, trust is all they had. The commitment was \emph{unverifiable}. No pasted signature, no diff, nothing the partner could check without access to the branch.


\vspace{0.35em}
\begin{tcolorbox}[failuregroup,colback=integ,
  title={\emojiicon{figs/emoji/1f9e9.png}\ \textbf{Commitment}: Failure to follow through on promises}]
  \causefigure{figs/Commitment.pdf}
\end{tcolorbox}

\paragraph{Communication.} In the third example, Agent A asks a direct question, ``Which approach would you prefer?'' The response is silence. Without an answer, the coordination loop collapses. A needed a decision to proceed, and without one, both agents continued with potentially incompatible assumptions. Unlike expectation failures (where information exists but isn't integrated) or commitment failures (where promises aren't kept), this is a failure to even establish shared context.

\vspace{0.35em}
\begin{tcolorbox}[failuregroup,colback=comm,
  title={\emojiicon{figs/emoji/1f4e3.png}\ \textbf{Communication}: Breakdown in using language to coordinate}]
  \causefigure{figs/Communication.pdf}
\end{tcolorbox}


\medskip
The examples above reveal why coordination, rather than raw coding ability, is often the limiting factor. The common thread is \emph{partial observability}. Each agent acts while holding an uncertain model of its partner's state, edits, and commitments. A merge can be conflict-free yet still embed incompatible assumptions.

These causes manifest through the symptoms in Tab.~\ref{tab:failure_symptoms}. Expectation failures produce work overlap and silent overwrites, commitment failures lead to unverifiable claims and broken promises, and communication failures result in unresponsiveness and repetition.

These failures suggest current models lack reliable representations for (i) \emph{partner state} (what the other agent has actually changed), (ii) \emph{checkable commitments} (contracts verifiable after merge), and (iii) \emph{cross-branch integration reasoning} (anticipating how independent patches interact). Coordination requires more than plausible code. It requires \emph{verifiable} and \emph{actionable} constraints for a partner operating under isolation. This explains why prompt optimization yields only marginal improvements (App.~\ref{app:prompt_optimization}). Most errors stem from coordination challenges, not prompt wording.

\noindent\textbf{The trust paradox.} We hypothesize that a deeper tension underlies expectation failures. Models are trained to be cautious, requiring observable evidence and resisting unverifiable assertions. This is a sensible default for single-agent interactions, where users may attempt to mislead the model. However, collaboration under workspace isolation requires the opposite. Agents must trust partner claims about states they cannot observe. When Agent~A reports ``I added the handler at line~50,'' Agent~B's instinct is to verify, but verification fails because they are on separate branches. This mismatch between \emph{verification-first} training and \emph{trust-requiring} collaboration may partly explain why agents consistently fail to update their model of partner state despite explicit communication.


Effective collaboration likely requires lightweight mechanisms that turn conversation into verifiable shared state, such as pasted signatures, explicit insertion-point contracts, and integration checks before declaring safety. We now turn to successful cases to see what these mechanisms look like in practice.

\subsection{Emergent Coordination Behavior}
\label{sec:emergent-coordination-behavior}
Among successful runs, we observe coordination patterns that are largely absent from failures. These behaviors are not prompted or scaffolded; they emerge when agents successfully navigate partial observability. What they share is a shift from vague intentions to specific commitments that a partner can verify even without seeing the underlying work. We identify three such patterns.
\vspace{-1em}
\paragraph{Role division.}
Agents agree on who handles which part of the task and establish clear boundaries around their scope.

\begin{figure}[htbp]
    \centering
    \includegraphics[width=\linewidth]{figs/emergent/Role_Division.pdf}
\end{figure}

What distinguishes successful role division is mutual confirmation.
Under partial observability, a unilateral declaration can easily be missed or misunderstood.
When both agents explicitly acknowledge the split, they create verified shared understanding that both sides can rely on during independent work.
\vspace{-1em}
\paragraph{Resource division.}
Agents avoid collisions by partitioning shared resources, most commonly specific files, code ranges, or \textit{ownership blocks}.

\begin{figure}[htbp]
    \centering
    \includegraphics[width=\linewidth]{figs/emergent/Resource_Division.pdf}
\end{figure}

What makes resource division effective is specificity.
Vague commitments cannot be verified and thus require trust.
Line-level boundaries, by contrast, create safe zones where conflict is impossible by construction.
\vspace{-1em}
\paragraph{Negotiation.}
Agents resolve conflicting approaches by proposing alternatives and converging on a single plan before acting.

\begin{figure}[htbp]
    \centering
    \includegraphics[width=\linewidth]{figs/emergent/Debate.pdf}
\end{figure}

Effective negotiation does cognitive work for both parties.
By proposing mutually exclusive options that fully specify what each agent will do, one agent reduces a complex coordination problem to a simple choice.
The result is not just agreement on intent but complete action specifications that leave nothing to interpret.

These coordination patterns are rare in our traces but their presence in successful cases suggests that the underlying capability exists.
The challenge is not teaching agents new coordination skills but making existing ones reliable.



